\import{function.tex}
\import{set/powerset.tex}
\import{set/equinumerosity.tex}

\section{Elementary function spaces}

\begin{proposition}\label{constant_graph_type}
    $A\times\{c\}\in\funs{A}{\{c\}}$.
\end{proposition}
\begin{proof}
    $A\times\{c\}$ is a relation by \auto{relation,times_elem_is_tuple}.
    $A\times\{c\}$ is right-unique
        by \auto{rightunique,times_tuple_elim,singleton_iff}.
    $\dom{A\times\{c\}}=A$
        by
        \auto{setext,dom_iff,times_tuple_intro,times_tuple_elim,singleton_iff}.
    For every $x\in\dom{A\times\{c\}}$ we have
        $(A\times\{c\})(x)\in\{c\}$
        by \auto{function_apply_elim,times_tuple_elim}.
    Follows by \obv{funs_intro}.
\end{proof}

\begin{proposition}\label{constant_graph_application}
    Let $x\in A$.
    Then $(A\times\{c\})(x)=c$.
\end{proposition}
\begin{proof}
    $(x,c)\in A\times\{c\}$
        by \auto{times_tuple_intro,singleton_iff}.
    $A\times\{c\}$ is a function by
        \auto{constant_graph_type,funs_is_function}.
    Follows by \obv{function_apply_intro}.
\end{proof}

\begin{proposition}\label{empty_domain_function_space}
    $\funs{\emptyset}{A}=\{\emptyset\}$.
\end{proposition}
\begin{proof}
    Show that every element of $\funs{\emptyset}{A}$ is an element of
        $\{\emptyset\}$.
    \begin{subproof}
        Fix $f\in\funs{\emptyset}{A}$.
        Hence $f$ is a function and $\dom{f}=\emptyset$
            by \obv{funs_is_function,funs}.
        Show that every element of $f$ is an element of $\emptyset$.
        \begin{subproof}
            Fix $w\in f$.
            Take $x,y$ such that $w=(x,y)$ by \obv{relation}.
            Hence $x\in\dom{f}$ by \auto{dom_iff}.
            Contradiction by \obv{emptyset}.
        \end{subproof}
        Hence $f\subseteq\emptyset$ by \obv{subseteq}.
        Hence $f=\emptyset$ by \obv{subseteq_emptyset_iff}.
        Follows by \obv{singleton_iff}.
    \end{subproof}
    Show that every element of $\{\emptyset\}$ is an element of
        $\funs{\emptyset}{A}$.
    \begin{subproof}
        Fix $f\in\{\emptyset\}$.
        Hence $f=\emptyset$ by \obv{singleton_iff}.
        $\emptyset$ is a function to $A$
            by \obv{codom_of_emptyset_can_be_anything}.
        $\dom{\emptyset}=\emptyset$ by \obv{dom_emptyset}.
        Hence $\emptyset\in\funs{\emptyset}{A}$ by \obv{funs_intro}.
    \end{subproof}
    Follows by set extensionality.
\end{proof}

\begin{proposition}\label{singleton_codom_function_space}
    $\funs{A}{\{c\}}=\{A\times\{c\}\}$.
\end{proposition}
\begin{proof}
    $A\times\{c\}\in\funs{A}{\{c\}}$
        by \obv{constant_graph_type}.
    Show that every element of $\funs{A}{\{c\}}$ is equal to
        $A\times\{c\}$.
    \begin{subproof}
        Fix $f\in\funs{A}{\{c\}}$.
        $f$ and $A\times\{c\}$ are functions and
            $\dom{f}=A=\dom{A\times\{c\}}$
            by
            \obv{funs_is_function,funs,constant_graph_type}.
        Show that for every $x$ we have $f(x)=(A\times\{c\})(x)$.
        \begin{subproof}
            Fix $x$.
            \begin{byCase}
                \caseOf{$x\in A$.}
                    $f(x)\in\{c\}$ by \obv{funs_type_apply}.
                    Hence $f(x)=c$ by \obv{singleton_iff}.
                    $(A\times\{c\})(x)=c$
                        by \obv{constant_graph_application}.
                \caseOf{$x\notin A$.}
                    $f(x)=\emptyset$ and
                        $(A\times\{c\})(x)=\emptyset$
                        by \obv{function_apply_default}.
            \end{byCase}
        \end{subproof}
        Follows by \obv{funext}.
    \end{subproof}
    Follows by \auto{setext,singleton_iff}.
\end{proof}

\begin{proposition}\label{singleton_domain_function_space_equinum}
    $\funs{\{a\}}{B}\approx B$.
\end{proposition}
\begin{proof}
    Let $e=\{(f,f(a))\mid f\in\funs{\{a\}}{B}\}$.
    $e$ is a relation by \auto{relation}.
    $e$ is right-unique by \auto{rightunique,pair_eq_iff}.
    $\dom{e}=\funs{\{a\}}{B}$
        by \auto{setext,dom_iff,pair_eq_iff}.
    For every $f\in\funs{\{a\}}{B}$ we have $e(f)\in B$
        by \auto{function_apply_intro,funs_type_apply,singleton_iff}.
    Hence $e\in\funs{\funs{\{a\}}{B}}{B}$ by \obv{funs_intro}.
    Show that $e$ is injective.
    \begin{subproof}
        Show that for all $f,g$ such that $f\in\dom{e}$ and
            $g\in\dom{e}$ and $e(f)=e(g)$ we have $f=g$.
        \begin{subproof}
            Fix $f,g$.
            Assume $f\in\dom{e}$ and $g\in\dom{e}$ and $e(f)=e(g)$.
            Hence $f,g\in\funs{\{a\}}{B}$.
            Hence $f$ and $g$ are functions and
                $\dom{f}=\{a\}=\dom{g}$
                by \obv{funs_is_function,funs}.
            $e(f)=f(a)$ and $e(g)=g(a)$
                by \auto{function_apply_intro}.
            Hence $f(a)=g(a)$.
            Show that for every $x$ we have $f(x)=g(x)$.
            \begin{subproof}
                Fix $x$.
                \begin{byCase}
                    \caseOf{$x\in\{a\}$.}
                        Hence $x=a$ by \obv{singleton_iff}.
                    \caseOf{$x\notin\{a\}$.}
                        $f(x)=\emptyset$ and $g(x)=\emptyset$
                            by \obv{function_apply_default}.
                \end{byCase}
            \end{subproof}
            Follows by \obv{funext}.
        \end{subproof}
        Follows by \auto{injective_function}.
    \end{subproof}
    Show that for every $b\in B$ there exists
        $f\in\funs{\{a\}}{B}$ such that $e(f)=b$.
    \begin{subproof}
        Fix $b\in B$.
        $\{a\}\times\{b\}\in\funs{\{a\}}{\{b\}}$
            by \obv{constant_graph_type}.
        Hence $\{a\}\times\{b\}\in\funs{\{a\}}{B}$
            by \obv{singleton_subset_intro,funs_weaken_codom}.
        $e(\{a\}\times\{b\})=(\{a\}\times\{b\})(a)$
            by \obv{function_apply_intro}.
        $(\{a\}\times\{b\})(a)=b$
            by \auto{constant_graph_application,singleton_iff}.
    \end{subproof}
    Hence $e\in\surj{\funs{\{a\}}{B}}{B}$ by \obv{surj}.
    Hence $e\in\bijections{\funs{\{a\}}{B}}{B}$
        by \obv{bijections}.
    Follows by \obv{equinum}.
\end{proof}

\begin{definition}\label{characteristic_graph}
    $\characteristicGraph{A}{S}{z}{o}
        =(S\times\{o\})\union((A\setminus S)\times\{z\})$.
\end{definition}

\begin{proposition}\label{characteristic_graph_type}
    Suppose $S\subseteq A$ and $z\neq o$.
    Then $\characteristicGraph{A}{S}{z}{o}\in\funs{A}{\{z,o\}}$.
\end{proposition}
\begin{proof}
    Let $f=\characteristicGraph{A}{S}{z}{o}$.
    $S\times\{o\}$ and $(A\setminus S)\times\{z\}$ are relations
        by \auto{relation,times_elem_is_tuple}.
    Hence $f$ is a relation
        by \auto{characteristic_graph,union_relations_is_relation}.
    Show that $f$ is right-unique.
    \begin{subproof}
        Show that for all $x,u,v$ such that $(x,u),(x,v)\in f$
            we have $u=v$.
        \begin{subproof}
            Fix $x,u,v$.
            Assume $(x,u),(x,v)\in f$.
            $(x,u)\in S\times\{o\}$ or
                $(x,u)\in(A\setminus S)\times\{z\}$
                by \auto{characteristic_graph,union_iff}.
            $(x,v)\in S\times\{o\}$ or
                $(x,v)\in(A\setminus S)\times\{z\}$
                by \auto{characteristic_graph,union_iff}.
            Follows by
                \auto{times_tuple_elim,singleton_iff,setminus}.
        \end{subproof}
        Follows by \obv{rightunique}.
    \end{subproof}
    $\dom{S\times\{o\}}=S$ and
        $\dom{(A\setminus S)\times\{z\}}=A\setminus S$
        by
        \auto{setext,dom_iff,times_tuple_intro,times_tuple_elim,singleton_iff}.
    Hence $\dom{f}=S\union(A\setminus S)$
        by \auto{characteristic_graph,dom_union}.
    Hence $\dom{f}=A$ by \obv{setminus_partition}.
    Show that for every $x\in\dom{f}$ we have $f(x)\in\{z,o\}$.
    \begin{subproof}
        Fix $x\in\dom{f}$.
        $(x,f(x))\in f$ by \obv{function_apply_elim}.
        $(x,f(x))\in S\times\{o\}$ or
            $(x,f(x))\in(A\setminus S)\times\{z\}$
            by \auto{characteristic_graph,union_iff}.
        Follows by \auto{times_tuple_elim,singleton_iff,cons_iff}.
    \end{subproof}
    Follows by \auto{funs_intro}.
\end{proof}

\begin{proposition}\label{characteristic_graph_one}
    Suppose $S\subseteq A$ and $z\neq o$.
    Let $x\in S$.
    Then $\apply{\characteristicGraph{A}{S}{z}{o}}{x}=o$.
\end{proposition}
\begin{proof}
    $(x,o)\in S\times\{o\}$
        by \auto{times_tuple_intro,singleton_iff}.
    Hence $(x,o)\in\characteristicGraph{A}{S}{z}{o}$
        by \auto{characteristic_graph,union_iff}.
    $\characteristicGraph{A}{S}{z}{o}$ is a function
        by \auto{characteristic_graph_type,funs_is_function}.
    Follows by \obv{function_apply_intro}.
\end{proof}

\begin{proposition}\label{characteristic_graph_zero}
    Suppose $S\subseteq A$ and $z\neq o$.
    Let $x\in A\setminus S$.
    Then $\apply{\characteristicGraph{A}{S}{z}{o}}{x}=z$.
\end{proposition}
\begin{proof}
    $(x,z)\in(A\setminus S)\times\{z\}$
        by \auto{times_tuple_intro,singleton_iff}.
    Hence $(x,z)\in\characteristicGraph{A}{S}{z}{o}$
        by \auto{characteristic_graph,union_iff}.
    $\characteristicGraph{A}{S}{z}{o}$ is a function
        by \auto{characteristic_graph_type,funs_is_function}.
    Follows by \obv{function_apply_intro}.
\end{proof}

\begin{theorem}\label{powerset_equinum_binary_function_space}
    Suppose $z\neq o$.
    Then $\pow{A}\approx\funs{A}{\{z,o\}}$.
\end{theorem}
\begin{proof}
    Let $e=\{(S,\characteristicGraph{A}{S}{z}{o})\mid S\in\pow{A}\}$.
    $e$ is a relation by \auto{relation}.
    $e$ is right-unique by \auto{rightunique,pair_eq_iff}.
    $\dom{e}=\pow{A}$ by \auto{setext,dom_iff,pair_eq_iff}.
    For every $S\in\pow{A}$ we have
        $e(S)=\characteristicGraph{A}{S}{z}{o}$
        by \obv{function_apply_intro}.
    For every $S\in\pow{A}$ we have
        $e(S)\in\funs{A}{\{z,o\}}$
        by \auto{powerset_elim,characteristic_graph_type}.
    Hence $e\in\funs{\pow{A}}{\funs{A}{\{z,o\}}}$
        by \obv{funs_intro}.
    Show that $e$ is injective.
    \begin{subproof}
        Show that for all $S,T$ such that $S\in\dom{e}$ and
            $T\in\dom{e}$ and $e(S)=e(T)$ we have $S=T$.
        \begin{subproof}
            Fix $S,T$.
            Assume $S\in\dom{e}$ and $T\in\dom{e}$ and $e(S)=e(T)$.
            Hence $S,T\in\pow{A}$.
            Hence $S\subseteq A$ and $T\subseteq A$
                by \obv{powerset_elim}.
            $e(S)=\characteristicGraph{A}{S}{z}{o}$ and
                $e(T)=\characteristicGraph{A}{T}{z}{o}$
                by \obv{function_apply_intro}.
            Show that every element of $S$ is an element of $T$.
            \begin{subproof}
                Fix $x\in S$.
                Suppose not.
                Hence $x\in A\setminus T$ by \obv{subseteq,setminus}.
                $\apply{\characteristicGraph{A}{S}{z}{o}}{x}=o$
                    by \obv{characteristic_graph_one}.
                $\apply{\characteristicGraph{A}{T}{z}{o}}{x}=z$
                    by \obv{characteristic_graph_zero}.
                Contradiction.
            \end{subproof}
            Hence $S\subseteq T$ by \obv{subseteq}.
            Show that every element of $T$ is an element of $S$.
            \begin{subproof}
                Fix $x\in T$.
                Suppose not.
                Hence $x\in A\setminus S$ by \obv{subseteq,setminus}.
                $\apply{\characteristicGraph{A}{T}{z}{o}}{x}=o$
                    by \obv{characteristic_graph_one}.
                $\apply{\characteristicGraph{A}{S}{z}{o}}{x}=z$
                    by \obv{characteristic_graph_zero}.
                Contradiction.
            \end{subproof}
            Hence $T\subseteq S$ by \obv{subseteq}.
            Follows by \obv{subseteq_antisymmetric}.
        \end{subproof}
        Follows by \auto{injective_function}.
    \end{subproof}
    Show that for every $f\in\funs{A}{\{z,o\}}$ there exists
        $S\in\pow{A}$ such that $e(S)=f$.
    \begin{subproof}
        Fix $f\in\funs{A}{\{z,o\}}$.
        Let $S=\{x\in A\mid f(x)=o\}$.
        $S\subseteq A$ by \obv{subseteq}.
        Hence $S\in\pow{A}$ by \obv{powerset_intro}.
        Let $g=\characteristicGraph{A}{S}{z}{o}$.
        $f$ and $g$ are functions and $\dom{f}=A=\dom{g}$
            by
            \auto{funs_is_function,funs,characteristic_graph_type}.
        Show that for every $x$ we have $f(x)=g(x)$.
        \begin{subproof}
            Fix $x$.
            \begin{byCase}
                \caseOf{$x\in A$.}
                    $f(x)\in\{z,o\}$ by \obv{funs_type_apply}.
                    \begin{byCase}
                        \caseOf{$f(x)=o$.}
                            Hence $x\in S$.
                            $g(x)=o$ by \auto{characteristic_graph_one}.
                        \caseOf{$f(x)\neq o$.}
                            Hence $f(x)=z$ by
                                \auto{cons_iff,singleton_iff}.
                            $x\in A\setminus S$ by \obv{setminus}.
                            $g(x)=z$ by \auto{characteristic_graph_zero}.
                    \end{byCase}
                \caseOf{$x\notin A$.}
                    $f(x)=\emptyset$ and $g(x)=\emptyset$
                        by \obv{function_apply_default}.
            \end{byCase}
        \end{subproof}
        Hence $f=g$ by \obv{funext}.
        $e(S)=g$ by \obv{function_apply_intro}.
    \end{subproof}
    Hence $e\in\surj{\pow{A}}{\funs{A}{\{z,o\}}}$ by \obv{surj}.
    Hence $e\in\bijections{\pow{A}}{\funs{A}{\{z,o\}}}$
        by \obv{bijections}.
    Follows by \obv{equinum}.
\end{proof}

\begin{theorem}\label{function_space_codom_injection}
    Let $f\in\inj{B}{D}$.
    Then there exists $H\in\inj{\funs{A}{B}}{\funs{A}{D}}$.
\end{theorem}
\begin{proof}
    $f\in\funs{B}{D}$ by \obv{inj}.
    Hence $f$ is a function and $\dom{f}=B$
        by \obv{funs_is_function,funs}.
    $f$ is injective by \auto{inj,injective_function}.
    Let $H=\{(h,f\circ h)\mid h\in\funs{A}{B}\}$.
    $H$ is a relation by \auto{relation}.
    $H$ is right-unique by \auto{rightunique,pair_eq_iff}.
    $\dom{H}=\funs{A}{B}$ by \auto{setext,dom_iff,pair_eq_iff}.
    For every $h\in\funs{A}{B}$ we have $f\circ h\in\funs{A}{D}$
        by \obv{funs_circ}.
    For every $h\in\funs{A}{B}$ we have $H(h)=f\circ h$
        by \obv{function_apply_intro}.
    Hence $H\in\funs{\funs{A}{B}}{\funs{A}{D}}$
        by \auto{funs_intro}.
    Show that $H$ is injective.
    \begin{subproof}
        Show that for all $h,k$ such that $h\in\dom{H}$ and
            $k\in\dom{H}$ and $H(h)=H(k)$ we have $h=k$.
        \begin{subproof}
            Fix $h,k$.
            Assume $h\in\dom{H}$ and $k\in\dom{H}$ and $H(h)=H(k)$.
            Hence $h,k\in\funs{A}{B}$.
            Hence $h$ and $k$ are functions and $\dom{h}=A=\dom{k}$
                by \obv{funs_is_function,funs}.
            $H(h)=f\circ h$ and $H(k)=f\circ k$
                by \obv{function_apply_intro}.
            Hence $f\circ h=f\circ k$.
            Show that for every $x$ we have $h(x)=k(x)$.
            \begin{subproof}
                Fix $x$.
                \begin{byCase}
                    \caseOf{$x\in A$.}
                        $h(x),k(x)\in B$ by \obv{funs_type_apply}.
                        $\ran{h}\subseteq B$ and $\ran{k}\subseteq B$
                            by \obv{funs_ran}.
                        Hence $\ran{h}\subseteq\dom{f}$ and
                            $\ran{k}\subseteq\dom{f}$.
                        $(f\circ h)(x)=f(h(x))$ and
                            $(f\circ k)(x)=f(k(x))$
                            by \obv{circ_apply}.
                        Hence $f(h(x))=f(k(x))$.
                        Follows by \obv{injective_function}.
                    \caseOf{$x\notin A$.}
                        $h(x)=\emptyset$ and $k(x)=\emptyset$
                            by \auto{function_apply_default}.
                \end{byCase}
            \end{subproof}
            Follows by \obv{funext}.
        \end{subproof}
        Follows by \auto{injective_function}.
    \end{subproof}
    Hence $H\in\inj{\funs{A}{B}}{\funs{A}{D}}$
        by \auto{inj,injective_function}.
\end{proof}

\subsection{Changing the domain of a function space}

An injection from $A$ into $C$ lets us extend a function on $A$ to $C$ once
a default codomain value has been specified.  Keeping the default explicit
isolates the exact nonemptiness requirement used by exponent monotonicity.

\begin{definition}\label{function_domain_extension}
    $\functionDomainExtension{e}{h}{C}{b}
        =\{(e(a),h(a))\mid a\in\dom{e}\}
            \union((C\setminus\ran{e})\times\{b\})$.
\end{definition}

\begin{proposition}\label{function_domain_extension_type}
    Let $e\in\inj{A}{C}$.
    Let $h\in\funs{A}{B}$.
    Let $b\in B$.
    Then $\functionDomainExtension{e}{h}{C}{b}\in\funs{C}{B}$.
\end{proposition}
\begin{proof}
    $e\in\funs{A}{C}$ by \obv{inj}.
    Hence $e$ is a function by \obv{funs_is_function}.
    $\dom{e}=A$ by \obv{funs}.
    Hence $e$ is injective by \auto{inj,injective_function}.
    $\dom{h}=A$ by \obv{funs}.
    Let $p=\{(e(a),h(a))\mid a\in\dom{e}\}$.
    Let $q=(C\setminus\ran{e})\times\{b\}$.
    Let $g=\functionDomainExtension{e}{h}{C}{b}$.
    $g=p\union q$ by \obv{function_domain_extension}.
    $p$ and $q$ are relations by \auto{relation,times_elem_is_tuple}.
    Hence $g$ is a relation by \obv{union_relations_is_relation}.
    Show that $g$ is right-unique.
    \begin{subproof}
        Show that for all $x,y,z$ such that $(x,y),(x,z)\in g$
            we have $y=z$.
        \begin{subproof}
            Fix $x,y,z$.
            Assume $(x,y),(x,z)\in g$.
            $(x,y)\in p$ or $(x,y)\in q$ by \auto{union_iff}.
            $(x,z)\in p$ or $(x,z)\in q$ by \auto{union_iff}.
            \begin{byCase}
                \caseOf{$(x,y)\in p$ and $(x,z)\in p$.}
                    Take $a,a'\in\dom{e}$ such that
                        $(x,y)=(e(a),h(a))$ and
                        $(x,z)=(e(a'),h(a'))$ by \auto.
                    $e(a)=e(a')$ and $y=h(a)$ and $z=h(a')$
                        by \auto{pair_eq_iff}.
                    Hence $a=a'$ by \obv{injective_function}.
                \caseOf{$(x,y)\in p$ and $(x,z)\in q$.}
                    Take $a\in\dom{e}$ such that
                        $(x,y)=(e(a),h(a))$.
                    $x=e(a)$ by \auto{pair_eq_iff}.
                    Hence $x\in\ran{e}$ by \auto{fun_ran_iff}.
                    $x\in C\setminus\ran{e}$
                        by \auto{times_tuple_elim}.
                    Contradiction by \obv{setminus}.
                \caseOf{$(x,y)\in q$ and $(x,z)\in p$.}
                    Take $a\in\dom{e}$ such that
                        $(x,z)=(e(a),h(a))$.
                    $x=e(a)$ by \auto{pair_eq_iff}.
                    Hence $x\in\ran{e}$ by \auto{fun_ran_iff}.
                    $x\in C\setminus\ran{e}$
                        by \auto{times_tuple_elim}.
                    Contradiction by \obv{setminus}.
                \caseOf{$(x,y)\in q$ and $(x,z)\in q$.}
                    $y=b$ and $z=b$
                        by \auto{times_tuple_elim,singleton_iff}.
            \end{byCase}
        \end{subproof}
        Follows by \obv{rightunique}.
    \end{subproof}
    $\dom{p}=\ran{e}$
        by \auto{setext,dom_iff,fun_ran_iff,pair_eq_iff}.
    $\dom{q}=C\setminus\ran{e}$
        by
        \auto{setext,dom_iff,times_tuple_intro,times_tuple_elim,singleton_iff}.
    Hence $\dom{g}=\ran{e}\union(C\setminus\ran{e})$
        by \auto{dom_union}.
    $\ran{e}\subseteq C$ by \obv{funs_ran}.
    Hence $\dom{g}=C$ by \obv{setminus_partition}.
    Show that for every $x\in\dom{g}$ we have $g(x)\in B$.
    \begin{subproof}
        Fix $x\in\dom{g}$.
        $(x,g(x))\in g$ by \obv{function_apply_elim}.
        $(x,g(x))\in p$ or $(x,g(x))\in q$ by \auto{union_iff}.
        \begin{byCase}
            \caseOf{$(x,g(x))\in p$.}
                Take $a\in\dom{e}$ such that
                    $(x,g(x))=(e(a),h(a))$.
                Hence $a\in A$.
                $g(x)=h(a)$ by \auto{pair_eq_iff}.
                $h(a)\in B$ by \obv{funs_type_apply}.
            \caseOf{$(x,g(x))\in q$.}
                $g(x)=b$ by \auto{times_tuple_elim,singleton_iff}.
        \end{byCase}
    \end{subproof}
    Hence $g$ is a function to $B$.
    Follows by \auto{funs_intro}.
\end{proof}

\begin{proposition}\label{function_domain_extension_on_range}
    Let $e\in\inj{A}{C}$.
    Let $h\in\funs{A}{B}$.
    Let $b\in B$.
    Let $a\in A$.
    Then $(\functionDomainExtension{e}{h}{C}{b})(e(a))=h(a)$.
\end{proposition}
\begin{proof}
    $e$ is a function and $\dom{e}=A$ by \auto{inj,funs_is_function,funs}.
    Hence $(e(a),h(a))\in\functionDomainExtension{e}{h}{C}{b}$
        by \auto{function_domain_extension,union_iff}.
    $\functionDomainExtension{e}{h}{C}{b}$ is a function
        by \auto{function_domain_extension_type,funs_is_function}.
    Follows by \obv{function_apply_intro}.
\end{proof}

\begin{proposition}\label{function_domain_extension_off_range}
    Let $e\in\inj{A}{C}$.
    Let $h\in\funs{A}{B}$.
    Let $b\in B$.
    Let $x\in C\setminus\ran{e}$.
    Then $(\functionDomainExtension{e}{h}{C}{b})(x)=b$.
\end{proposition}
\begin{proof}
    $(x,b)\in(C\setminus\ran{e})\times\{b\}$
        by \auto{times_tuple_intro,singleton_iff}.
    Hence $(x,b)\in\functionDomainExtension{e}{h}{C}{b}$
        by \auto{function_domain_extension,union_iff}.
    $\functionDomainExtension{e}{h}{C}{b}$ is a function
        by \auto{function_domain_extension_type,funs_is_function}.
    Follows by \obv{function_apply_intro}.
\end{proof}

\begin{theorem}\label{function_space_domain_injection_with_default}
    Let $e\in\inj{A}{C}$.
    Let $b\in B$.
    Then there exists $H\in\inj{\funs{A}{B}}{\funs{C}{B}}$.
\end{theorem}
\begin{proof}
    Let $H=\{(h,\functionDomainExtension{e}{h}{C}{b})
        \mid h\in\funs{A}{B}\}$.
    $H$ is a relation by \auto{relation}.
    $H$ is right-unique by \auto{rightunique,pair_eq_iff}.
    $\dom{H}=\funs{A}{B}$ by \auto{setext,dom_iff,pair_eq_iff}.
    For every $h\in\funs{A}{B}$ we have
        $\functionDomainExtension{e}{h}{C}{b}\in\funs{C}{B}$
        by \obv{function_domain_extension_type}.
    For every $h\in\funs{A}{B}$ we have
        $H(h)=\functionDomainExtension{e}{h}{C}{b}$
        by \obv{function_apply_intro}.
    Hence $H\in\funs{\funs{A}{B}}{\funs{C}{B}}$
        by \auto{funs_intro}.
    Show that $H$ is injective.
    \begin{subproof}
        Show that for all $h,k$ such that $h\in\dom{H}$ and
            $k\in\dom{H}$ and $H(h)=H(k)$ we have $h=k$.
        \begin{subproof}
            Fix $h,k$.
            Assume $h\in\dom{H}$ and $k\in\dom{H}$ and $H(h)=H(k)$.
            Hence $h,k\in\funs{A}{B}$.
            Hence $h$ and $k$ are functions and $\dom{h}=A=\dom{k}$
                by \obv{funs_is_function,funs}.
            $H(h)=\functionDomainExtension{e}{h}{C}{b}$ and
                $H(k)=\functionDomainExtension{e}{k}{C}{b}$
                by \obv{function_apply_intro}.
            Hence $\functionDomainExtension{e}{h}{C}{b}
                =\functionDomainExtension{e}{k}{C}{b}$.
            Show that for every $a$ we have $h(a)=k(a)$.
            \begin{subproof}
                Fix $a$.
                \begin{byCase}
                    \caseOf{$a\in A$.}
                        $(\functionDomainExtension{e}{h}{C}{b})(e(a))
                            =h(a)$ and
                            $(\functionDomainExtension{e}{k}{C}{b})(e(a))
                            =k(a)$
                            by \obv{function_domain_extension_on_range}.
                        \begin{align*}
                            h(a)
                                &= (\functionDomainExtension{e}{h}{C}{b})(e(a))
                                \\
                                &= (\functionDomainExtension{e}{k}{C}{b})(e(a))
                                \\
                                &= k(a)
                        \end{align*}
                    \caseOf{$a\notin A$.}
                        $h(a)=\emptyset$ and $k(a)=\emptyset$
                            by \auto{function_apply_default}.
                \end{byCase}
            \end{subproof}
            Follows by \obv{funext}.
        \end{subproof}
        Follows by \auto{injective_function}.
    \end{subproof}
    Hence $H\in\inj{\funs{A}{B}}{\funs{C}{B}}$
        by \auto{inj,injective_function}.
\end{proof}

\begin{corollary}\label{function_space_domain_injection}
    Let $e\in\inj{A}{C}$.
    Suppose $B$ is inhabited.
    Then there exists $H\in\inj{\funs{A}{B}}{\funs{C}{B}}$.
\end{corollary}
\begin{proof}
    Take $b$ such that $b\in B$.
    Follows by \obv{function_space_domain_injection_with_default}.
\end{proof}


\subsection{Currying}

\begin{definition}\label{function_section}
    $\functionSection{f}{a}{B}
        =\{(b,f((a,b)))\mid b\in B\}$.
\end{definition}

\begin{proposition}\label{function_section_type}
    Let $f\in\funs{A\times B}{C}$.
    Let $a\in A$.
    Then $\functionSection{f}{a}{B}\in\funs{B}{C}$.
\end{proposition}
\begin{proof}
    Let $s=\functionSection{f}{a}{B}$.
    $s$ is a relation by \auto{function_section,relation}.
    $s$ is right-unique
        by \auto{function_section,rightunique,pair_eq_iff}.
    $\dom{s}=B$ by \auto{function_section,setext,dom_iff,pair_eq_iff}.
    For every $b\in B$ we have $(a,b)\in A\times B$
        by \obv{times_tuple_intro}.
    Hence for every $b\in B$ we have $f((a,b))\in C$
        by \obv{funs_type_apply}.
    For every $b\in B$ we have $s(b)=f((a,b))$
        by \obv{function_section,function_apply_intro}.
    Hence $s$ is a function to $C$ by \auto.
    Follows by \auto{funs_intro}.
\end{proof}

\begin{proposition}\label{function_section_application}
    Let $f\in\funs{A\times B}{C}$.
    Let $a\in A$.
    Let $b\in B$.
    Then $(\functionSection{f}{a}{B})(b)=f((a,b))$.
\end{proposition}
\begin{proof}
    $(b,f((a,b)))\in\functionSection{f}{a}{B}$
        by \obv{function_section}.
    $\functionSection{f}{a}{B}$ is a function
        by \auto{function_section_type,funs_is_function}.
    Follows by \obv{function_apply_intro}.
\end{proof}

\begin{definition}\label{function_curry}
    $\functionCurry{f}{A}{B}
        =\{(a,\functionSection{f}{a}{B})\mid a\in A\}$.
\end{definition}

\begin{proposition}\label{function_curry_type}
    Let $f\in\funs{A\times B}{C}$.
    Then $\functionCurry{f}{A}{B}\in\funs{A}{\funs{B}{C}}$.
\end{proposition}
\begin{proof}
    Let $g=\functionCurry{f}{A}{B}$.
    $g$ is a relation by \auto{function_curry,relation}.
    $g$ is right-unique
        by \auto{function_curry,rightunique,pair_eq_iff}.
    $\dom{g}=A$ by \auto{function_curry,setext,dom_iff,pair_eq_iff}.
    For every $a\in A$ we have
        $\functionSection{f}{a}{B}\in\funs{B}{C}$
        by \obv{function_section_type}.
    For every $a\in A$ we have
        $g(a)=\functionSection{f}{a}{B}$
        by \obv{function_curry,function_apply_intro}.
    Hence $g$ is a function to $\funs{B}{C}$ by \auto.
    Follows by \auto{funs_intro}.
\end{proof}

\begin{proposition}\label{function_curry_application}
    Let $f\in\funs{A\times B}{C}$.
    Let $a\in A$.
    Then $(\functionCurry{f}{A}{B})(a)=\functionSection{f}{a}{B}$.
\end{proposition}
\begin{proof}
    $(a,\functionSection{f}{a}{B})\in\functionCurry{f}{A}{B}$
        by \obv{function_curry}.
    $\functionCurry{f}{A}{B}$ is a function
        by \auto{function_curry_type,funs_is_function}.
    Follows by \obv{function_apply_intro}.
\end{proof}

\begin{definition}\label{function_uncurry}
    $\functionUncurry{g}{A}{B}
        =\{((a,b),(g(a))(b))\mid a\in A,b\in B\}$.
\end{definition}

\begin{proposition}\label{function_uncurry_type}
    Let $g\in\funs{A}{\funs{B}{C}}$.
    Then $\functionUncurry{g}{A}{B}\in\funs{A\times B}{C}$.
\end{proposition}
\begin{proof}
    Let $f=\functionUncurry{g}{A}{B}$.
    $f$ is a relation by \auto{function_uncurry,relation}.
    $f$ is right-unique
        by \auto{function_uncurry,rightunique,pair_eq_iff}.
    Show that every element of $\dom{f}$ is an element of $A\times B$.
    \begin{subproof}
        Fix $x\in\dom{f}$.
        Take $y$ such that $(x,y)\in f$ by \obv{dom_iff}.
        Take $a,b$ such that $a\in A$ and $b\in B$ and
            $(x,y)=((a,b),(g(a))(b))$ by \auto{function_uncurry}.
        Hence $x=(a,b)$ by \auto{pair_eq_iff}.
        Follows by \auto{times_tuple_intro}.
    \end{subproof}
    Hence $\dom{f}\subseteq A\times B$ by \obv{subseteq}.
    Show that every element of $A\times B$ is an element of $\dom{f}$.
    \begin{subproof}
        Fix $x\in A\times B$.
        Take $a,b$ such that $x=(a,b)$ by \auto{times_elem_is_tuple}.
        Hence $a\in A$ and $b\in B$ by \obv{times_tuple_elim}.
        $((a,b),(g(a))(b))\in f$ by \obv{function_uncurry}.
        Follows by \auto{dom_iff}.
    \end{subproof}
    Hence $A\times B\subseteq\dom{f}$ by \obv{subseteq}.
    Hence $\dom{f}=A\times B$ by \obv{setext}.
    Show that for every $x\in\dom{f}$ we have $f(x)\in C$.
    \begin{subproof}
        Fix $x\in\dom{f}$.
        Hence $x\in A\times B$.
        Take $a,b$ such that $x=(a,b)$ by \auto{times_elem_is_tuple}.
        Hence $a\in A$ and $b\in B$ by \obv{times_tuple_elim}.
        $g(a)\in\funs{B}{C}$ by \obv{funs_type_apply}.
        Hence $(g(a))(b)\in C$ by \obv{funs_type_apply}.
        $f((a,b))=(g(a))(b)$
            by \auto{function_uncurry,function_apply_intro}.
    \end{subproof}
    Hence $f$ is a function to $C$.
    Follows by \auto{funs_intro}.
\end{proof}

\begin{proposition}\label{function_uncurry_application}
    Let $g\in\funs{A}{\funs{B}{C}}$.
    Let $a\in A$.
    Let $b\in B$.
    Then $(\functionUncurry{g}{A}{B})((a,b))=(g(a))(b)$.
\end{proposition}
\begin{proof}
    $((a,b),(g(a))(b))\in\functionUncurry{g}{A}{B}$
        by \obv{function_uncurry}.
    $\functionUncurry{g}{A}{B}$ is a function
        by \auto{function_uncurry_type,funs_is_function}.
    Follows by \obv{function_apply_intro}.
\end{proof}

\begin{proposition}\label{function_uncurry_curry}
    Let $f\in\funs{A\times B}{C}$.
    Then $\functionUncurry{\functionCurry{f}{A}{B}}{A}{B}=f$.
\end{proposition}
\begin{proof}
    Let $g=\functionCurry{f}{A}{B}$.
    Let $u=\functionUncurry{g}{A}{B}$.
    $g\in\funs{A}{\funs{B}{C}}$ by \auto{function_curry_type}.
    Hence $u\in\funs{A\times B}{C}$ by \auto{function_uncurry_type}.
    Hence $u$ and $f$ are functions by \obv{funs_is_function}.
    $\dom{u}=A\times B=\dom{f}$ by \obv{funs}.
    Show that for every $x$ we have $u(x)=f(x)$.
    \begin{subproof}
        Fix $x$.
        \begin{byCase}
            \caseOf{$x\in A\times B$.}
                Take $a,b$ such that $x=(a,b)$
                    by \auto{times_elem_is_tuple}.
                Hence $a\in A$ and $b\in B$ by \obv{times_tuple_elim}.
                $u((a,b))=(g(a))(b)$
                    by \auto{function_uncurry_application}.
                $g(a)=\functionSection{f}{a}{B}$
                    by \auto{function_curry_application}.
                $(\functionSection{f}{a}{B})(b)=f((a,b))$
                    by \obv{function_section_application}.
            \caseOf{$x\notin A\times B$.}
                $x\notin\dom{u}$ and $x\notin\dom{f}$.
                $u(x)=\emptyset$ and $f(x)=\emptyset$
                    by \obv{function_apply_default}.
        \end{byCase}
    \end{subproof}
    Hence $u=f$ by \obv{funext}.
\end{proof}

\begin{proposition}\label{function_curry_uncurry}
    Let $g\in\funs{A}{\funs{B}{C}}$.
    Then $\functionCurry{\functionUncurry{g}{A}{B}}{A}{B}=g$.
\end{proposition}
\begin{proof}
    Let $f=\functionUncurry{g}{A}{B}$.
    Let $h=\functionCurry{f}{A}{B}$.
    $f\in\funs{A\times B}{C}$ by \auto{function_uncurry_type}.
    Hence $h\in\funs{A}{\funs{B}{C}}$ by \auto{function_curry_type}.
    Hence $h$ and $g$ are functions by \obv{funs_is_function}.
    $\dom{h}=A=\dom{g}$ by \obv{funs}.
    Show that for every $a$ we have $h(a)=g(a)$.
    \begin{subproof}
        Fix $a$.
        \begin{byCase}
            \caseOf{$a\in A$.}
                $h(a)=\functionSection{f}{a}{B}$
                    by \auto{function_curry_application}.
                $\functionSection{f}{a}{B}\in\funs{B}{C}$
                    by \obv{function_section_type}.
                $g(a)\in\funs{B}{C}$ by \obv{funs_type_apply}.
                Hence $\functionSection{f}{a}{B}$ and $g(a)$ are functions
                    by \obv{funs_is_function}.
                $\dom{\functionSection{f}{a}{B}}=B=\dom{g(a)}$
                    by \obv{funs}.
                Show that for every $b$ we have
                    $(\functionSection{f}{a}{B})(b)=(g(a))(b)$.
                \begin{subproof}
                    Fix $b$.
                    \begin{byCase}
                        \caseOf{$b\in B$.}
                            $(\functionSection{f}{a}{B})(b)=f((a,b))$
                                by \obv{function_section_application}.
                            $f((a,b))=(g(a))(b)$
                                by \auto{function_uncurry_application}.
                        \caseOf{$b\notin B$.}
                            $b\notin\dom{\functionSection{f}{a}{B}}$ and
                                $b\notin\dom{g(a)}$.
                            $(\functionSection{f}{a}{B})(b)=\emptyset$ and
                                $(g(a))(b)=\emptyset$
                                by \obv{function_apply_default}.
                    \end{byCase}
                \end{subproof}
                Hence $\functionSection{f}{a}{B}=g(a)$
                    by \obv{funext}.
            \caseOf{$a\notin A$.}
                $a\notin\dom{h}$ and $a\notin\dom{g}$.
                $h(a)=\emptyset$ and $g(a)=\emptyset$
                    by \obv{function_apply_default}.
        \end{byCase}
    \end{subproof}
    Hence $h=g$ by \obv{funext}.
\end{proof}

\begin{theorem}\label{function_space_currying_equinumerous}
    $\funs{A\times B}{C}\approx\funs{A}{\funs{B}{C}}$.
\end{theorem}
\begin{proof}
    Let $H=\{(f,\functionCurry{f}{A}{B})
        \mid f\in\funs{A\times B}{C}\}$.
    $H$ is a relation by \auto{relation}.
    $H$ is right-unique by \auto{rightunique,pair_eq_iff}.
    $\dom{H}=\funs{A\times B}{C}$
        by \auto{setext,dom_iff,pair_eq_iff}.
    For every $f\in\funs{A\times B}{C}$ we have
        $\functionCurry{f}{A}{B}\in\funs{A}{\funs{B}{C}}$
        by \obv{function_curry_type}.
    For every $f\in\funs{A\times B}{C}$ we have
        $H(f)=\functionCurry{f}{A}{B}$
        by \obv{function_apply_intro}.
    Hence $H\in\funs{\funs{A\times B}{C}}{\funs{A}{\funs{B}{C}}}$
        by \auto{funs_intro}.
    Show that $H$ is injective.
    \begin{subproof}
        Show that for all $f,k$ such that $f\in\dom{H}$ and
            $k\in\dom{H}$ and $H(f)=H(k)$ we have $f=k$.
        \begin{subproof}
            Fix $f,k$.
            Assume $f\in\dom{H}$ and $k\in\dom{H}$ and $H(f)=H(k)$.
            Hence $f,k\in\funs{A\times B}{C}$.
            $H(f)=\functionCurry{f}{A}{B}$ and
                $H(k)=\functionCurry{k}{A}{B}$
                by \obv{function_apply_intro}.
            Hence $\functionCurry{f}{A}{B}=\functionCurry{k}{A}{B}$.
            $\functionUncurry{\functionCurry{f}{A}{B}}{A}{B}=f$ and
                $\functionUncurry{\functionCurry{k}{A}{B}}{A}{B}=k$
                by \obv{function_uncurry_curry}.
        \end{subproof}
        Follows by \auto{injective_function}.
    \end{subproof}
    Show that for every $g\in\funs{A}{\funs{B}{C}}$ there exists
        $f\in\funs{A\times B}{C}$ such that $H(f)=g$.
    \begin{subproof}
        Fix $g\in\funs{A}{\funs{B}{C}}$.
        Let $f=\functionUncurry{g}{A}{B}$.
        $f\in\funs{A\times B}{C}$ by \auto{function_uncurry_type}.
        Hence $H(f)=\functionCurry{f}{A}{B}$
            by \obv{function_apply_intro}.
        $\functionCurry{f}{A}{B}=g$ by \auto{function_curry_uncurry}.
    \end{subproof}
    Hence $H\in\surj{\funs{A\times B}{C}}{\funs{A}{\funs{B}{C}}}$
        by \obv{surj}.
    Hence $H\in\bijections{\funs{A\times B}{C}}{\funs{A}{\funs{B}{C}}}$
        by \obv{bijections}.
    Follows by \obv{equinum}.
\end{proof}
